% `tex.web` §1048 lists `\eqno` among the forbidden cases — `non_math(eq_no)`
% — and §1140 admits it in math mode only when `privileged`:
%
%     mmode+eq_no: if privileged then
%       if cur_group=math_shift_group then start_eq_no
%       else off_save;
%
%     function privileged:boolean;
%     begin if mode>0 then privileged:=true
%     else begin report_illegal_case; privileged:=false; end;
%     end;
%
% `mode>0` is the display: `$…$` pushes `-mmode`, so an equation number inside
% ordinary math is §1049's illegal case just as one in a paragraph is.  Four
% modes, four messages, and §1049 names each of them:
%
%     print_err("You can't use `"); print_cmd_chr(cur_cmd,cur_chr);
%     print("' in "); print_mode(mode);
%
% This engine admitted `\eqno` silently everywhere it did not apply, so
% trip.tex line 254's bare `\eqno` — which the file's own comment calls
% "another error" — produced none.
\catcode`\{=1 \catcode`\}=2 \catcode`\$=3
\font\r=cmr10 \r \hsize100pt \parindent0pt
A\eqno B\par
\vbox{\eqno}
$ x\eqno y$\par
\leqno
\end
